% generated by tools/make_tex_tables.py -- do not edit
% The paragraph that belongs to tables/table10_codegap.tex: what the columns mean and how they were
% read. It is body prose, input by the section that owns the table, and not a second
% block of small type inside the float.
Table~\ref{tab:codegap} is the whole of the paper-against-code finding, in the form the finding is made: a public repository, the commit \texttt{tools/fetch\_code.py} cloned, the file inside it, and the two values. \emph{file in it, and where} gives the path as that repository spells it and, beneath it, the function, configuration key or branch the value sits in. \emph{what the paper prints} names the table or equation the value was read from. \emph{what that file contains} is what is in the parsed copy under \path{code/md}, which is the same snapshot every other claim in this survey about that repository is made from. The commit is the one in \texttt{corpus/code\_manifest.json}, printed to ten characters; the rows print the comparison, while the reproducibility materials retain each row's full text, confidence and review note. \emph{conf.} is the row's recorded confidence: 8 high, 1 medium. The medium-confidence row is PenSpin, pending a direct repository read; the reason is recorded in the accompanying materials. No cell states a cause, and none is a claim about what the work's authors did: a reader with a browser can check every line of this table without asking anyone.
